% --- Archon physics formalization source begin ---
% archon:physics
% archon:covers PhyXMiniProblems/problem_phyx_mini_0707.lean
% archon:source-report reports/phyx_mini/problem_phyx_mini_0707.source.json

\chapter{Physics problem phyx\_mini\_0707}
\label{ch:PhyXMiniProblems_problem_phyx_mini_0707}

\paragraph{Problem source.}
\#\# Physical scenario

Figure shows the forces acting on our simplified model of the forearm. The biceps pulls the forearm up against the upper arm at the elbow, so the force \textbackslash{}(\textbackslash{}vec\{F\}\_\{\textbackslash{}mathrm\{elbow\}\}\textbackslash{}) on the forearm at the elbow—a force due to the upper arm—is a downward force.

\#\# Question

What is the magnitude of \textbackslash{}(\textbackslash{}vec\{F\}\_\{\textbackslash{}mathrm\{elbow\}\}\textbackslash{})?

\#\# Answer choices

- A: A: 3800N

- B: B: 4000N

- C: C: 3900N

- D: D: 3700N

\#\# Figure caption (auxiliary; use the image as primary evidence)

The figure illustrates a biomechanical setup involving forces and distances related to arm movement and loading at the elbow joint. 

1. **Objects and Labels:**
   - A vertical gray bar represents the upper arm.
   - A horizontal pink bar represents the forearm.
   - Three arrows indicate forces: 
     - \textbackslash{}( \textbackslash{}vec\{F\}\_\{\textbackslash{}text\{tendon\}\} \textbackslash{}) is an upward force located between the elbow and the hand.
     - \textbackslash{}( \textbackslash{}vec\{F\}\_\{\textbackslash{}text\{elbow\}\} \textbackslash{}) is a downward force at the elbow joint.
     - \textbackslash{}( \textbackslash{}vec\{F\}\_\{\textbackslash{}text\{barbell\}\} \textbackslash{}) is a downward force at the end of the forearm representing the barbell weight.
   - Distances are labeled: 
     - \textbackslash{}( d\_\{\textbackslash{}text\{tendon\}\} \textbackslash{}) is the distance from the elbow to the tendon force application point.
     - \textbackslash{}( d\_\{\textbackslash{}text\{arm\}\} \textbackslash{}) is the distance from the elbow to where the barbell force is applied.

2. **Text and Formulas:**
   - Known quantities:
     - \textbackslash{}( d\_\{\textbackslash{}text\{tendon\}\} = 4.0 \textbackslash{}, \textbackslash{}text\{cm\} \textbackslash{})
     - \textbackslash{}( d\_\{\textbackslash{}text\{arm\}\} = 35 \textbackslash{}, \textbackslash{}text\{cm\} \textbackslash{})
     - \textbackslash{}( F\_\{\textbackslash{}text\{barbell\}\} = 450 \textbackslash{}, \textbackslash{}text\{N\} \textbackslash{})
   - Objective: Find \textbackslash{}( F\_\{\textbackslash{}text\{tendon\}\} \textbackslash{}).
   - A blue dashed arrow indicates the direction of torque and states, "These forces cause torques about the elbow."

3. **Relationships:**
   - The diagram visually shows how forces and distances interact to produce torques around the elbow joint.
   - It sets up the problem of solving for the tendon force based on the given distances and known force (barbell).

\#\# Dataset metadata

Statics; Physical Model Grounding Reasoning, Spatial Relation Reasoning

\paragraph{Recorded answer/context.}
C: C: 3900N

\paragraph{Figure/image path.}
/root/proposal\_for\_physic/science-mango/phyx\_mini\_run/phyx\_data/test\_image/707.png

\paragraph{Formalization target.}
create a compiling Lean file with sorry bodies at `PhyXMiniProblems/problem\_phyx\_mini\_0707.lean`. The Lean declarations must preserve the physical quantities, dimensions or dimensional roles, figure labels, governing-law hypotheses, and final relation expressed by this problem.
Use Mathlib/Physlib names found through LeanExplore where available. If a physics API is missing, introduce faithful local abstractions rather than scalar placeholder aliases.

\begin{theorem}[Physics formalization target]
\label{thm:physics:phyx_mini_0707:target}
\lean{PhyXMiniProblems.ProblemPhyXMini0707.problem_phyx_mini_0707}
\uses{def:physics:phyx-mini-0707:phyxminiproblems-problemphyxmini0707-forcemagnitudeinnewtons, def:physics:phyx-mini-0707:phyxminiproblems-problemphyxmini0707-forearmstaticssetup, def:physics:phyx-mini-0707:phyxminiproblems-problemphyxmini0707-matchessuppliedfigure, def:physics:phyx-mini-0707:phyxminiproblems-problemphyxmini0707-matchesproblemreadouts, def:physics:phyx-mini-0707:phyxminiproblems-problemphyxmini0707-matchesforearmgeometry, def:physics:phyx-mini-0707:phyxminiproblems-problemphyxmini0707-matchesdisplayedforcedirections, def:physics:phyx-mini-0707:phyxminiproblems-problemphyxmini0707-satisfiesforearmstaticequilibrium, def:physics:phyx-mini-0707:phyxminiproblems-problemphyxmini0707-recordedanswerchoice, def:physics:phyx-mini-0707:phyxminiproblems-problemphyxmini0707-isuniquenearestdisplayedforce}
The assigned autoformalize agent should translate this physics problem into Lean declarations in the covered file, with theorem and lemma proof bodies written as `by sorry`.
\end{theorem}
\begin{proof}
This is an autoformalization task, not a proof task. Produce statements that can later be proved by the physics prover without weakening the physical model.
\end{proof}

% --- Archon declaration topology begin ---
\section{Declaration topology}
\begin{definition}[force Dimension]
  \label{def:physics:phyx-mini-0707:phyxminiproblems-problemphyxmini0707-forcedimension}
  \lean{PhyXMiniProblems.ProblemPhyXMini0707.forceDimension}
  The physical dimension M L T⁻² of force
\end{definition}
\begin{proof}
  This is the declared model definition of force Dimension.
\end{proof}
\begin{definition}[torque Dimension]
  \label{def:physics:phyx-mini-0707:phyxminiproblems-problemphyxmini0707-torquedimension}
  \lean{PhyXMiniProblems.ProblemPhyXMini0707.torqueDimension}
  \uses{def:physics:phyx-mini-0707:phyxminiproblems-problemphyxmini0707-forcedimension}
  The physical dimension M L² T⁻² of torque
\end{definition}
\begin{proof}
  This is the declared model definition of torque Dimension.
\end{proof}
\begin{definition}[Length Quantity]
  \label{def:physics:phyx-mini-0707:phyxminiproblems-problemphyxmini0707-lengthquantity}
  \lean{PhyXMiniProblems.ProblemPhyXMini0707.LengthQuantity}
  A unit-independent nonnegative physical length
\end{definition}
\begin{proof}
  This is the declared model definition of Length Quantity.
\end{proof}
\begin{definition}[Planar Position Quantity]
  \label{def:physics:phyx-mini-0707:phyxminiproblems-problemphyxmini0707-planarpositionquantity}
  \lean{PhyXMiniProblems.ProblemPhyXMini0707.PlanarPositionQuantity}
  A unit-independent signed position vector in the figure's plane
\end{definition}
\begin{proof}
  This is the declared model definition of Planar Position Quantity.
\end{proof}
\begin{definition}[Planar Force Quantity]
  \label{def:physics:phyx-mini-0707:phyxminiproblems-problemphyxmini0707-planarforcequantity}
  \lean{PhyXMiniProblems.ProblemPhyXMini0707.PlanarForceQuantity}
  \uses{def:physics:phyx-mini-0707:phyxminiproblems-problemphyxmini0707-forcedimension}
  A unit-independent physical force vector in the figure's plane
\end{definition}
\begin{proof}
  This is the declared model definition of Planar Force Quantity.
\end{proof}
\begin{definition}[Planar Torque Quantity]
  \label{def:physics:phyx-mini-0707:phyxminiproblems-problemphyxmini0707-planartorquequantity}
  \lean{PhyXMiniProblems.ProblemPhyXMini0707.PlanarTorqueQuantity}
  \uses{def:physics:phyx-mini-0707:phyxminiproblems-problemphyxmini0707-torquedimension}
  The dimensionful type occupied by a signed planar torque component
\end{definition}
\begin{proof}
  This is the declared model definition of Planar Torque Quantity.
\end{proof}
\begin{definition}[x Axis]
  \label{def:physics:phyx-mini-0707:phyxminiproblems-problemphyxmini0707-xaxis}
  \lean{PhyXMiniProblems.ProblemPhyXMini0707.xAxis}
  Coordinate 0, the horizontal axis pointing right in the figure
\end{definition}
\begin{proof}
  This is the declared model definition of x Axis.
\end{proof}
\begin{definition}[y Axis]
  \label{def:physics:phyx-mini-0707:phyxminiproblems-problemphyxmini0707-yaxis}
  \lean{PhyXMiniProblems.ProblemPhyXMini0707.yAxis}
  Coordinate 1, the vertical axis pointing up in the figure
\end{definition}
\begin{proof}
  This is the declared model definition of y Axis.
\end{proof}
\begin{definition}[length In Meters]
  \label{def:physics:phyx-mini-0707:phyxminiproblems-problemphyxmini0707-lengthinmeters}
  \lean{PhyXMiniProblems.ProblemPhyXMini0707.lengthInMeters}
  \uses{def:physics:phyx-mini-0707:phyxminiproblems-problemphyxmini0707-lengthquantity}
  Coherent-SI metre readout of a nonnegative physical length
\end{definition}
\begin{proof}
  This is the declared model definition of length In Meters.
\end{proof}
\begin{definition}[length In Centimeters]
  \label{def:physics:phyx-mini-0707:phyxminiproblems-problemphyxmini0707-lengthincentimeters}
  \lean{PhyXMiniProblems.ProblemPhyXMini0707.lengthInCentimeters}
  \uses{def:physics:phyx-mini-0707:phyxminiproblems-problemphyxmini0707-lengthquantity, def:physics:phyx-mini-0707:phyxminiproblems-problemphyxmini0707-lengthinmeters}
  Centimetre readout used by the two distance labels in the image
\end{definition}
\begin{proof}
  This is the declared model definition of length In Centimeters.
\end{proof}
\begin{definition}[position In Meters]
  \label{def:physics:phyx-mini-0707:phyxminiproblems-problemphyxmini0707-positioninmeters}
  \lean{PhyXMiniProblems.ProblemPhyXMini0707.positionInMeters}
  \uses{def:physics:phyx-mini-0707:phyxminiproblems-problemphyxmini0707-planarpositionquantity}
  Coherent-SI position-vector readout, with coordinates in metres
\end{definition}
\begin{proof}
  This is the declared model definition of position In Meters.
\end{proof}
\begin{definition}[force Vector In Newtons]
  \label{def:physics:phyx-mini-0707:phyxminiproblems-problemphyxmini0707-forcevectorinnewtons}
  \lean{PhyXMiniProblems.ProblemPhyXMini0707.forceVectorInNewtons}
  \uses{def:physics:phyx-mini-0707:phyxminiproblems-problemphyxmini0707-planarforcequantity}
  Coherent-SI force-vector readout, with coordinates in newtons
\end{definition}
\begin{proof}
  This is the declared model definition of force Vector In Newtons.
\end{proof}
\begin{definition}[force Magnitude In Newtons]
  \label{def:physics:phyx-mini-0707:phyxminiproblems-problemphyxmini0707-forcemagnitudeinnewtons}
  \lean{PhyXMiniProblems.ProblemPhyXMini0707.forceMagnitudeInNewtons}
  \uses{def:physics:phyx-mini-0707:phyxminiproblems-problemphyxmini0707-planarforcequantity, def:physics:phyx-mini-0707:phyxminiproblems-problemphyxmini0707-forcevectorinnewtons}
  Euclidean magnitude, in newtons, of a physical planar force
\end{definition}
\begin{proof}
  This is the declared model definition of force Magnitude In Newtons.
\end{proof}
\begin{definition}[planar Torque In Newton Meters]
  \label{def:physics:phyx-mini-0707:phyxminiproblems-problemphyxmini0707-planartorqueinnewtonmeters}
  \lean{PhyXMiniProblems.ProblemPhyXMini0707.planarTorqueInNewtonMeters}
  \uses{def:physics:phyx-mini-0707:phyxminiproblems-problemphyxmini0707-planarpositionquantity, def:physics:phyx-mini-0707:phyxminiproblems-problemphyxmini0707-planarforcequantity, def:physics:phyx-mini-0707:phyxminiproblems-problemphyxmini0707-xaxis, def:physics:phyx-mini-0707:phyxminiproblems-problemphyxmini0707-yaxis, def:physics:phyx-mini-0707:phyxminiproblems-problemphyxmini0707-positioninmeters, def:physics:phyx-mini-0707:phyxminiproblems-problemphyxmini0707-forcevectorinnewtons}
  The out-of-plane component of (applicationPoint - pivot) × force, read in newton-metres. Its sign is positive for counterclockwise torque in the figure's right-handed coordinate convention
\end{definition}
\begin{proof}
  This is the declared model definition of planar Torque In Newton Meters.
\end{proof}
\begin{definition}[Applied Force Label]
  \label{def:physics:phyx-mini-0707:phyxminiproblems-problemphyxmini0707-appliedforcelabel}
  \lean{PhyXMiniProblems.ProblemPhyXMini0707.AppliedForceLabel}
  The three force arrows named in the problem and primary image
\end{definition}
\begin{proof}
  This is the declared model definition of Applied Force Label.
\end{proof}
\begin{definition}[Application Point Label]
  \label{def:physics:phyx-mini-0707:phyxminiproblems-problemphyxmini0707-applicationpointlabel}
  \lean{PhyXMiniProblems.ProblemPhyXMini0707.ApplicationPointLabel}
  Distinguished force-application locations along the forearm
\end{definition}
\begin{proof}
  This is the declared model definition of Application Point Label.
\end{proof}
\begin{definition}[Figure Element]
  \label{def:physics:phyx-mini-0707:phyxminiproblems-problemphyxmini0707-figureelement}
  \lean{PhyXMiniProblems.ProblemPhyXMini0707.FigureElement}
  Rigid-body elements visibly represented in the primary image
\end{definition}
\begin{proof}
  This is the declared model definition of Figure Element.
\end{proof}
\begin{definition}[Figure Quantity Label]
  \label{def:physics:phyx-mini-0707:phyxminiproblems-problemphyxmini0707-figurequantitylabel}
  \lean{PhyXMiniProblems.ProblemPhyXMini0707.FigureQuantityLabel}
  Literal physical-quantity labels printed in the supplied raster
\end{definition}
\begin{proof}
  This is the declared model definition of Figure Quantity Label.
\end{proof}
\begin{definition}[Arrow Direction]
  \label{def:physics:phyx-mini-0707:phyxminiproblems-problemphyxmini0707-arrowdirection}
  \lean{PhyXMiniProblems.ProblemPhyXMini0707.ArrowDirection}
  Qualitative directions used by the three force arrows
\end{definition}
\begin{proof}
  This is the declared model definition of Arrow Direction.
\end{proof}
\begin{definition}[expected Arrow Direction]
  \label{def:physics:phyx-mini-0707:phyxminiproblems-problemphyxmini0707-expectedarrowdirection}
  \lean{PhyXMiniProblems.ProblemPhyXMini0707.expectedArrowDirection}
  \uses{def:physics:phyx-mini-0707:phyxminiproblems-problemphyxmini0707-appliedforcelabel, def:physics:phyx-mini-0707:phyxminiproblems-problemphyxmini0707-arrowdirection}
  The direction visibly assigned to each force arrow
\end{definition}
\begin{proof}
  This is the declared model definition of expected Arrow Direction.
\end{proof}
\begin{definition}[expected Application Point]
  \label{def:physics:phyx-mini-0707:phyxminiproblems-problemphyxmini0707-expectedapplicationpoint}
  \lean{PhyXMiniProblems.ProblemPhyXMini0707.expectedApplicationPoint}
  \uses{def:physics:phyx-mini-0707:phyxminiproblems-problemphyxmini0707-appliedforcelabel, def:physics:phyx-mini-0707:phyxminiproblems-problemphyxmini0707-applicationpointlabel}
  The application point visibly assigned to each force arrow
\end{definition}
\begin{proof}
  This is the declared model definition of expected Application Point.
\end{proof}
\begin{definition}[Forearm Statics Figure]
  \label{def:physics:phyx-mini-0707:phyxminiproblems-problemphyxmini0707-forearmstaticsfigure}
  \lean{PhyXMiniProblems.ProblemPhyXMini0707.ForearmStaticsFigure}
  \uses{def:physics:phyx-mini-0707:phyxminiproblems-problemphyxmini0707-appliedforcelabel, def:physics:phyx-mini-0707:phyxminiproblems-problemphyxmini0707-applicationpointlabel, def:physics:phyx-mini-0707:phyxminiproblems-problemphyxmini0707-figureelement, def:physics:phyx-mini-0707:phyxminiproblems-problemphyxmini0707-figurequantitylabel, def:physics:phyx-mini-0707:phyxminiproblems-problemphyxmini0707-arrowdirection}
  Qualitative and literal content transcribed from image 707.png
\end{definition}
\begin{proof}
  This is the declared model definition of Forearm Statics Figure.
\end{proof}
\begin{definition}[Forearm Statics Setup]
  \label{def:physics:phyx-mini-0707:phyxminiproblems-problemphyxmini0707-forearmstaticssetup}
  \lean{PhyXMiniProblems.ProblemPhyXMini0707.ForearmStaticsSetup}
  \uses{def:physics:phyx-mini-0707:phyxminiproblems-problemphyxmini0707-lengthquantity, def:physics:phyx-mini-0707:phyxminiproblems-problemphyxmini0707-planarpositionquantity, def:physics:phyx-mini-0707:phyxminiproblems-problemphyxmini0707-planarforcequantity, def:physics:phyx-mini-0707:phyxminiproblems-problemphyxmini0707-appliedforcelabel, def:physics:phyx-mini-0707:phyxminiproblems-problemphyxmini0707-forearmstaticsfigure}
  Independent quantities in the simplified forearm experiment. No requested force magnitude or answer choice is defined into this setup
\end{definition}
\begin{proof}
  This is the declared model definition of Forearm Statics Setup.
\end{proof}
\begin{definition}[Matches Supplied Figure]
  \label{def:physics:phyx-mini-0707:phyxminiproblems-problemphyxmini0707-matchessuppliedfigure}
  \lean{PhyXMiniProblems.ProblemPhyXMini0707.MatchesSuppliedFigure}
  \uses{def:physics:phyx-mini-0707:phyxminiproblems-problemphyxmini0707-expectedarrowdirection, def:physics:phyx-mini-0707:phyxminiproblems-problemphyxmini0707-expectedapplicationpoint, def:physics:phyx-mini-0707:phyxminiproblems-problemphyxmini0707-forearmstaticssetup}
  Primary-image evidence: the two arm segments, all five quantity labels, all three arrows and their application points, and the torque annotation
\end{definition}
\begin{proof}
  This is the declared model definition of Matches Supplied Figure.
\end{proof}
\begin{definition}[Matches Problem Readouts]
  \label{def:physics:phyx-mini-0707:phyxminiproblems-problemphyxmini0707-matchesproblemreadouts}
  \lean{PhyXMiniProblems.ProblemPhyXMini0707.MatchesProblemReadouts}
  \uses{def:physics:phyx-mini-0707:phyxminiproblems-problemphyxmini0707-lengthincentimeters, def:physics:phyx-mini-0707:phyxminiproblems-problemphyxmini0707-forcemagnitudeinnewtons, def:physics:phyx-mini-0707:phyxminiproblems-problemphyxmini0707-forearmstaticssetup}
  The three numerical values printed in the figure. Neither the tendon-force magnitude nor the elbow-force magnitude occurs among these source readouts
\end{definition}
\begin{proof}
  This is the declared model definition of Matches Problem Readouts.
\end{proof}
\begin{definition}[Matches Forearm Geometry]
  \label{def:physics:phyx-mini-0707:phyxminiproblems-problemphyxmini0707-matchesforearmgeometry}
  \lean{PhyXMiniProblems.ProblemPhyXMini0707.MatchesForearmGeometry}
  \uses{def:physics:phyx-mini-0707:phyxminiproblems-problemphyxmini0707-xaxis, def:physics:phyx-mini-0707:phyxminiproblems-problemphyxmini0707-yaxis, def:physics:phyx-mini-0707:phyxminiproblems-problemphyxmini0707-lengthinmeters, def:physics:phyx-mini-0707:phyxminiproblems-problemphyxmini0707-positioninmeters, def:physics:phyx-mini-0707:phyxminiproblems-problemphyxmini0707-forearmstaticssetup}
  Coordinate realization of the horizontal forearm: the elbow is the origin, and both force application points lie to its right on the horizontal axis at the two named physical distances
\end{definition}
\begin{proof}
  This is the declared model definition of Matches Forearm Geometry.
\end{proof}
\begin{definition}[Matches Displayed Force Directions]
  \label{def:physics:phyx-mini-0707:phyxminiproblems-problemphyxmini0707-matchesdisplayedforcedirections}
  \lean{PhyXMiniProblems.ProblemPhyXMini0707.MatchesDisplayedForceDirections}
  \uses{def:physics:phyx-mini-0707:phyxminiproblems-problemphyxmini0707-xaxis, def:physics:phyx-mini-0707:phyxminiproblems-problemphyxmini0707-yaxis, def:physics:phyx-mini-0707:phyxminiproblems-problemphyxmini0707-forcevectorinnewtons, def:physics:phyx-mini-0707:phyxminiproblems-problemphyxmini0707-forcemagnitudeinnewtons, def:physics:phyx-mini-0707:phyxminiproblems-problemphyxmini0707-forearmstaticssetup}
  The force directions shown in the image, expressed as coherent-SI Cartesian components. These equations state only directions; no unknown force value is inserted
\end{definition}
\begin{proof}
  This is the declared model definition of Matches Displayed Force Directions.
\end{proof}
\begin{definition}[Satisfies Forearm Static Equilibrium]
  \label{def:physics:phyx-mini-0707:phyxminiproblems-problemphyxmini0707-satisfiesforearmstaticequilibrium}
  \lean{PhyXMiniProblems.ProblemPhyXMini0707.SatisfiesForearmStaticEquilibrium}
  \uses{def:physics:phyx-mini-0707:phyxminiproblems-problemphyxmini0707-forcevectorinnewtons, def:physics:phyx-mini-0707:phyxminiproblems-problemphyxmini0707-planartorqueinnewtonmeters, def:physics:phyx-mini-0707:phyxminiproblems-problemphyxmini0707-forearmstaticssetup}
  Static translational and rotational equilibrium of the forearm. This is the governing physical law and contains neither requested numerical magnitude nor an answer-choice label
\end{definition}
\begin{proof}
  This is the declared model definition of Satisfies Forearm Static Equilibrium.
\end{proof}
\begin{definition}[Answer Choice]
  \label{def:physics:phyx-mini-0707:phyxminiproblems-problemphyxmini0707-answerchoice}
  \lean{PhyXMiniProblems.ProblemPhyXMini0707.AnswerChoice}
  The four answer labels supplied with the prose question
\end{definition}
\begin{proof}
  This is the declared model definition of Answer Choice.
\end{proof}
\begin{definition}[displayed Force In Newtons]
  \label{def:physics:phyx-mini-0707:phyxminiproblems-problemphyxmini0707-displayedforceinnewtons}
  \lean{PhyXMiniProblems.ProblemPhyXMini0707.displayedForceInNewtons}
  \uses{def:physics:phyx-mini-0707:phyxminiproblems-problemphyxmini0707-answerchoice}
  Force magnitude in newtons printed beside each answer label
\end{definition}
\begin{proof}
  This is the declared model definition of displayed Force In Newtons.
\end{proof}
\begin{definition}[recorded Answer Choice]
  \label{def:physics:phyx-mini-0707:phyxminiproblems-problemphyxmini0707-recordedanswerchoice}
  \lean{PhyXMiniProblems.ProblemPhyXMini0707.recordedAnswerChoice}
  \uses{def:physics:phyx-mini-0707:phyxminiproblems-problemphyxmini0707-answerchoice}
  The answer label recorded by the supplied dataset metadata
\end{definition}
\begin{proof}
  This is the declared model definition of recorded Answer Choice.
\end{proof}
\begin{definition}[Is Nearest Displayed Force]
  \label{def:physics:phyx-mini-0707:phyxminiproblems-problemphyxmini0707-isnearestdisplayedforce}
  \lean{PhyXMiniProblems.ProblemPhyXMini0707.IsNearestDisplayedForce}
  \uses{def:physics:phyx-mini-0707:phyxminiproblems-problemphyxmini0707-planarforcequantity, def:physics:phyx-mini-0707:phyxminiproblems-problemphyxmini0707-forcemagnitudeinnewtons, def:physics:phyx-mini-0707:phyxminiproblems-problemphyxmini0707-answerchoice, def:physics:phyx-mini-0707:phyxminiproblems-problemphyxmini0707-displayedforceinnewtons}
  A choice is at least as close to a force magnitude as every alternative
\end{definition}
\begin{proof}
  This is the declared model definition of Is Nearest Displayed Force.
\end{proof}
\begin{definition}[Is Unique Nearest Displayed Force]
  \label{def:physics:phyx-mini-0707:phyxminiproblems-problemphyxmini0707-isuniquenearestdisplayedforce}
  \lean{PhyXMiniProblems.ProblemPhyXMini0707.IsUniqueNearestDisplayedForce}
  \uses{def:physics:phyx-mini-0707:phyxminiproblems-problemphyxmini0707-planarforcequantity, def:physics:phyx-mini-0707:phyxminiproblems-problemphyxmini0707-answerchoice, def:physics:phyx-mini-0707:phyxminiproblems-problemphyxmini0707-isnearestdisplayedforce}
  A choice is the unique nearest displayed force magnitude
\end{definition}
\begin{proof}
  This is the declared model definition of Is Unique Nearest Displayed Force.
\end{proof}
% --- Archon declaration topology end ---
% --- Archon physics formalization source end ---
